% U-CODE PROSPECTIVE MANUSCRIPT SHELL — NOT SUBMISSION READY
% Official Springer Nature class must be obtained from the official template:
% https://www.overleaf.com/latex/templates/springer-nature-latex-template/myxmhdsbzkyd
% Journal requirement: \documentclass{sn-jnl}
\documentclass[pdflatex,sn-basic]{sn-jnl}

\usepackage{amsmath,amssymb}
\usepackage{graphicx}
\usepackage{booktabs}

\begin{document}

\title[U-CODE]{U-CODE: Uncertainty-Guided Search for Reliable AI-Generated Software}

% Known intended corresponding author; coauthors/affiliations must be approved
% by the actual human research team before submission.
\author*[1]{\fnm{Chandan Kumar} \sur{Sah}}\email{sahchandan98@buaa.edu.cn}
\affil*[1]{\orgdiv{School of Computer Science and Engineering},
  \orgname{Beihang University},
  \orgaddress{\city{Beijing}, \country{China}}}

\abstract{%
% PLACEHOLDER: Replace with a factual 150--250-word abstract AFTER empirical
% results are complete. Do not invent numerical improvements or data sizes.
This manuscript shell defines a prospective research study of uncertainty-guided
search for AI-generated software. The research will investigate whether
calibrated pre-verification uncertainty can guide the allocation of finite
inference and software verification budgets across candidate generation,
retrieval, testing, static analysis, repair, and abstention. The proposed
evaluation will distinguish naturally occurring agent-authored patches from
new model generations on executable software engineering tasks. No empirical
results are claimed in this initial manuscript scaffold.
}

\keywords{AI-generated software, uncertainty calibration,
search-based software engineering, agentic code verification}

\maketitle

\section{Introduction}\label{sec:introduction}
% TODO: motivating instance with verified source and independent test.
% TODO: problem gap and contributions AFTER evidence and novelty audit.

\section{Related Work}\label{sec:related}
% TODO: 20--30 verified scholarly references and original contribution map.

\section{Problem Formulation}\label{sec:formulation}
Given observable state $s_t$, an agent chooses intervention $a_t$ under
budget $B$, seeking measurable improvements in independently verified
reliability, safety and resource use. The selected action may maximize a
predicted acquisition objective based on feasible multi-objective gain:
\begin{equation}
  a_t^\star =
  \arg\max_{a \in \mathcal A_t,\; c(a)\le B_t}
  \frac{\mathbb E[\Delta HV \mid s_t,a]}{c(a)}.
  \label{eq:ehvi}
\end{equation}
% This is a research hypothesis and example acquisition, not an established
% novel or empirically validated algorithm. Specify assumptions in final text.

\section{U-CODE Method}\label{sec:method}
% TODO: agent architecture, risk calibration, exploration/exploitation,
% policy ablations and exact search algorithm.

\section{Dataset Provenance and Experimental Design}\label{sec:dataset}
% TODO: screened naturalistic agent PR cohort; distinguish PR/file/line/task.
% TODO: model-generated benchmark cohort and independent evaluation oracles.
% TODO: sampling and rights checks, data splits, model revisions and budgets.

\section{Results}\label{sec:results}
% NO RESULTS YET. Fill only from executable, version-pinned trace manifests.
% TODO: RQ1--RQ6, effect estimates, CIs, calibration, risk--coverage,
% language/model heterogeneity, budget curves, comparisons and negative cases.

\section{Discussion and Threats to Validity}\label{sec:discussion}
% TODO: bot authorship uncertainty, selection bias, code/license restrictions,
% hidden test incompleteness, contamination and practical impact.

\section{Conclusions}\label{sec:conclusions}
% TODO: only conclusions actually supported by completed experiments.

\section*{Statements and Declarations}
\textbf{Funding.} To be completed and verified by the authors.

\textbf{Competing interests.} To be declared by the authors.

\textbf{Data and code availability.} Source-code and metadata releases are
subject to source licenses and data-sharing approvals. Final URLs and
reproducibility statements will be provided after artifact validation.

\textbf{Author contributions.} To be completed and approved by all authors.

\textbf{Use of AI tools.} Any non-copy-editing AI assistance in research or
manuscript preparation must be disclosed as required by the publisher.

% Add verified references to references.bib and uncomment ONLY when cited:
% \bibliography{references}

\end{document}
